\begin{frame}{클리핑의 비대칭: 상한은 ``개선 방향''에만}
  \small
  두 패널을 함께 읽으면 min의 역할이 정확히 보인다:
  \begin{itemize}
    \item \kb{범위 안}($r_t\in[0.8,1.2]$): 원래 중요도 목적함수
      $r_tA_t$ 그대로 — $r_t=1$ 근처(데이터 수집 직후)에서는
      REINFORCE 계열 그래디언트와 본질적으로 같은 방향.
    \item \kb{범위 밖, 개선 방향}: ``좋은 것을 더 밀어주는 것''과
      ``나쁜 것을 더 억제하는 것''의 \textbf{추가 이득이 0}
      (검은 곡선 수평).
    \item \kb{범위 밖, 악화 방향}: ``나쁜 것을 더 밀어올리는 것''
     에는 \textbf{상한이 없다}.
  \end{itemize}
  \vskip0.5em
  \begin{block}{미분으로 보면 더 명확}
    목적함수가 수평인 구간에서는 $r_t$ 에 대한 기울기가 정확히 0 —
    ``추가 이득이 0''은 문자 그대로 ``이 방향으로는 그만''.
    \kb{min이 ``비관적(더 작은)쪽''을 고르기 때문에} 자연스럽게
    생기는 비대칭.
  \end{block}
\end{frame}
